\section{Kruskal's Algorithm (Minimum Spanning Tree)}
\label{sec:graph-kruskal}

\problemstatement{Given a weighted, connected, undirected graph, find a
\textbf{Minimum Spanning Tree} (MST) -- a subset of edges connecting
all nodes, with no cycles, at the minimum possible total weight.}

\begin{patternbox}
\emph{``Connect everything as cheaply as possible''} is solved
greedily: sort all edges by weight, and consider them from
\textbf{cheapest to most expensive}, adding each edge to the MST
unless it would create a \textbf{cycle} -- checked instantly using
Section~\ref{sec:graph-disjoint-set-impl}'s Disjoint Set
($\textit{find}$ on both endpoints; if equal, they're already
connected, so this edge would close a cycle and is skipped).
\end{patternbox}

\subsection*{Why the greedy choice is correct (the cut property)}

For any partition of the nodes into two non-empty groups, the
\textbf{cheapest} edge crossing between them must belong to \emph{some}
MST -- if it didn't, swapping it in for whatever more expensive edge
\emph{does} cross between the groups in that MST would only decrease
the total weight, contradicting minimality. Processing edges from
cheapest to most expensive and greedily accepting any edge that
connects two still-separate components is exactly repeatedly applying
this cut property, which is why the greedy strategy provably produces
an MST.

\subsection*{Algorithm}

\begin{enumerate}
  \item Sort all edges by weight, ascending.
  \item Initialize a Disjoint Set over all $n$ nodes.
  \item For each edge $(u,v,w)$ in sorted order: if $\textit{find}(u)
        \neq \textit{find}(v)$, add the edge to the MST (accumulate
        its weight) and $\textit{unite}(u,v)$.
  \item Stop once $n{-}1$ edges have been added (a spanning tree on
        $n$ nodes always has exactly $n{-}1$ edges).
\end{enumerate}

\subsection*{C++ implementation}

\begin{lstlisting}
#include <bits/stdc++.h>
using namespace std;

class DisjointSet {
    vector<int> parent, size;
public:
    DisjointSet(int n) : parent(n), size(n, 1) {
        for (int i = 0; i < n; i++) parent[i] = i;
    }
    int find(int x) {
        if (parent[x] == x) return x;
        return parent[x] = find(parent[x]);
    }
    void unite(int x, int y) {
        int rootX = find(x), rootY = find(y);
        if (rootX == rootY) return;
        if (size[rootX] < size[rootY]) swap(rootX, rootY);
        parent[rootY] = rootX;
        size[rootX] += size[rootY];
    }
};

int kruskalMST(int n, vector<vector<int>>& edges) {
    sort(edges.begin(), edges.end(),
         [](const vector<int>& a, const vector<int>& b) { return a[2] < b[2]; });

    DisjointSet ds(n);
    int totalWeight = 0, edgesUsed = 0;

    for (auto& e : edges) {
        int u = e[0], v = e[1], w = e[2];
        if (ds.find(u) != ds.find(v)) {
            ds.unite(u, v);
            totalWeight += w;
            edgesUsed++;
            if (edgesUsed == n - 1) break;
        }
    }
    return totalWeight;
}

int main() {
    vector<vector<int>> edges = {{0,1,4},{0,2,1},{1,2,2},{1,3,5},{2,3,8}};
    cout << kruskalMST(4, edges) << endl; // 8
    return 0;
}
\end{lstlisting}

\subsection*{Dry run}

\dryrun{Take edges $0{-}1(4),0{-}2(1),1{-}2(2),1{-}3(5),2{-}3(8)$,
$n=4$.}

Sorted: $0{-}2(1),1{-}2(2),0{-}1(4),1{-}3(5),2{-}3(8)$. Take $0{-}2$
(weight $1$): different components, add; total $=1$. Take $1{-}2$
(weight $2$): $\textit{find}(1) \neq \textit{find}(2)$, add; total
$=3$; now $\{0,1,2\}$ connected. Take $0{-}1$ (weight $4$):
$\textit{find}(0)=\textit{find}(1)$ (both already in $\{0,1,2\}$'s
group) -- \textbf{skip}, would form a cycle. Take $1{-}3$ (weight
$5$): different components ($3$ is still separate), add; total
$=8$. Now $3$ edges used ($0{-}2$, $1{-}2$, $1{-}3$) $=n{-}1=3$, stop.
Final MST weight: $8$ -- node $3$'s cheapest available connection in
this graph is the weight-$5$ edge to node $1$, since the alternative
$2{-}3$ edge costs $8$.

\begin{complexitybox}
\textbf{Time Complexity.} $O(E \log E)$ for sorting, plus
$O(E \cdot \alpha(V))$ for the Disjoint Set operations:
$O(E \log E)$ overall.
\textbf{Space Complexity.} $O(V)$ for the Disjoint Set, plus $O(E)$ for
sorting the edge list.
\end{complexitybox}

\begin{takeawaybox}
Kruskal's algorithm is a direct, clean composition of two ideas already
established elsewhere in this book series: sorting to consider choices
in a greedy-optimal order (the same idea behind every Greedy chapter
algorithm), and Disjoint Set to reject cycle-forming choices in
near-constant time. Neither idea is new here -- their combination is.
\end{takeawaybox}
